% Volume IV: Law as an Epistemic Machine
% Part VIII. Legal Friction as Freedom
% Chapter 41: The Inconvenience Principle
% Spine: proof-spines/ch041-spine.md

\chapter{The Inconvenience Principle}
\label{ch:ch041}

The \textbf{inconvenience principle} states the central proposition of the monograph: where information is incomplete, actors are fallible, and coercive actions are asymmetrically costly, procedural inconvenience can enlarge substantive freedom. Its technical handle is a friction parameter \(\phi\), which delays intervention just enough to force justification before power hardens into fact. Social loss can then be modeled by
\[
L(\phi)
=
L_{\mathrm{false\ action}}(\phi)
+
L_{\mathrm{missed\ action}}(\phi)
+
L_{\mathrm{delay}}(\phi).
\]
\ToyModel\ Too little friction produces impulsive coercion; too much blocks legitimate protection.
The task of legal design is therefore not to maximize friction, but to locate it where epistemic uncertainty meets irreversible power. Friction changes sign and becomes attrition when the waiting it imposes does no evidentiary work, falls asymmetrically on those already exposed to power, and offers no clear route to completion.
